% Where --disable-resolvability-validation moves the failure. Referenced from
% chapter 10. Bare tikzpicture; the figure environment, caption and label live
% at the call site.
%
% Two columns for the same fault, and one horizontal rule. Everything above the
% rule happens on a laptop with both schemas open; everything below it happens
% in production. The columns run top to bottom rather than left to right
% because the right-hand one has three outcomes and no room for them sideways.
\begin{tikzpicture}[
    font=\small,
    stage/.style={draw, rounded corners=2pt, minimum width=46mm,
                  minimum height=8mm, font=\scriptsize, align=center},
    result/.style={draw, rounded corners=2pt, minimum width=52mm,
                   minimum height=12mm, font=\scriptsize, align=center,
                   fill=black!8},
    quiet/.style={draw, rounded corners=2pt, minimum width=52mm,
                  minimum height=11mm, font=\scriptsize, align=center,
                  dashed},
    flow/.style={-{Stealth[length=2mm]}},
    lbl/.style={font=\scriptsize, align=center, inner sep=1pt},
    divider/.style={gray, dotted, thick}
  ]

  % -- the two columns -------------------------------------------------------
  \node[lbl] at (3.5,0.5) {\textbf{Composed as it stands}};
  \node[lbl] at (10.5,0.5) {\textbf{Composed with the check off}};

  \node[stage] (cmdA) at (3.5,-0.4)  {\code{wgc router compose}};
  \node[stage] (cmdB) at (10.5,-0.4)
    {\code{wgc router compose}\\\code{--disable-resolvability-validation}};

  \node[result] (errA) at (3.5,-2.0)
    {4 errors, exit 1\\each names a field\\and the path a client takes to it};
  \node[result] (cfgB) at (10.5,-2.0)
    {a config, exit 0\\one property differs\\from a healthy one};

  \draw[flow] (cmdA) -- (errA);
  \draw[flow] (cmdB) -- (cfgB);

  % -- the divide ------------------------------------------------------------
  \draw[divider] (0.3,-3.1) -- (13.7,-3.1);
  \node[lbl, gray, anchor=south east] at (13.7,-3.03) {build time};
  \node[lbl, gray, anchor=north east] at (13.7,-3.17) {request time};

  % -- what each one leaves you with -----------------------------------------
  \node[lbl, align=center] (noneA) at (3.5,-4.1)
    {nothing is deployed,\\because nothing\\was produced};

  \node[quiet] (routerB) at (10.5,-4.1)
    {the router starts\\no warning about the graph};
  \draw[flow] (cfgB) -- (routerB);

  \node[result] (okB) at (10.5,-5.8)
    {a query inside one subgraph\\HTTP 200};
  \node[result] (badB) at (10.5,-7.4)
    {a query across the boundary\\HTTP 500,
     \code{internal server error}};

  % Two outcomes of the same router, not a sequence, so the second one is
  % routed around the first rather than through it.
  \draw[flow] (routerB.south) -- (okB.north);
  \draw[flow] (routerB.west) -- (7.4,-4.1) -- (7.4,-7.4) -- (badB.west);

\end{tikzpicture}
